% ---------------------------------------------------------------------------
% Resultados da RSL 1 (OSINT para órgãos policiais e de inteligência).
% Fontes: docs/triagem/prisma-etapa2.md, docs/triagem/snowball-leitura/,
% docs/triagem/qualidade/, docs/triagem/sintese/ (mapeamento e leituras A1 a A4).
% ---------------------------------------------------------------------------
\section{RSL 1: ferramentas OSINT para órgãos policiais e de inteligência}
\label{sec:resultados-rsl1}

Esta seção apresenta a seleção, as características e a qualidade dos estudos
incluídos na RSL 1, a síntese temática de cada questão de pesquisa e a
ligação dos achados com o levantamento de requisitos do buscador. Os estudos
incluídos são citados pelo identificador entre parênteses, por exemplo
(RSL1-0494). O prefixo RSL1 indica estudos da busca nas bases, e o prefixo
SB1, estudos do \textit{snowballing}. A lista completa, com a referência de
cada relato, está no Apêndice~\ref{ap:estudos}.

A leitura do texto completo na etapa 2, o \textit{snowballing}, a avaliação
de qualidade, a extração e a síntese foram feitos pela IA como terceira
opinião, com regras de aplicação aceitas pelo autor (R1). O orientador (R2)
ainda confirma essas regras, e a avaliação independente de R1 e R2 ainda não
foi feita. Os resultados desta seção são, portanto, provisórios. A síntese
codificou os campos da extração, de até 60 palavras cada, e não o texto
completo. Uma ausência registrada aqui (por exemplo, ``nenhum estudo trata do
Brasil'') significa ausência nesses campos. Os números de resultado dos
estudos foram conferidos por \textit{script} contra o texto completo.
% PENDENTE: R2 confirma as regras R1-A a R1-D (etapa 2), R1-E a R1-H
% (snowballing), Q-A a Q-C (qualidade) e S-A a S-D (síntese). Se alguma regra
% mudar, os números e as listas de IDs desta seção são recalculados por script.

% ---------------------------------------------------------------------------
\subsection{Seleção dos estudos}
\label{sec:rsl1-selecao}

A Tabela~\ref{tab:prisma-rsl1} mostra o fluxo da busca nas bases. As cinco
fontes retornaram 1.928 registros. Depois da remoção de 789 duplicatas,
1.139 registros foram triados por título, resumo e palavras-chave. A etapa 1
excluiu 689, e 450 passaram para a etapa 2.
% PENDENTE: o registro de buscas da RSL 1 diz "1.923 registros" na seção de
% ampliação, mas a soma dos arquivos dá 1.928 (1.794 + 134). Conferir e
% corrigir o registro de buscas (docs/triagem/prisma-etapa2.md).

Na etapa 2, o recorte da Seção~\ref{sec:metodologia-selecao} (CI5, ferramenta
com avaliação; CI6, uso por órgão policial, de investigação criminal ou de
inteligência explícito no registro; CI7, publicação a partir de 2015) excluiu
361 registros pelo resumo e pelos metadados: 187 pelo CI6, 88 pelo CI5 e 86
pelo CI7. Outros 19 registros eram estudos secundários e ficaram para a
análise à parte. Dos 70 relatos buscados, 21 não foram recuperados pelo
acesso da UnB/CAPES nem em acesso aberto. Dos 49 relatos lidos, 26 foram
excluídos no texto completo, 13 deles pelo CI5. Entraram 23
estudos primários.

O \textit{snowballing} (Tabela~\ref{tab:snowball-rsl1}) partiu desses 23
estudos e parou na sexta iteração, a primeira sem estudo novo. Ele
identificou 2.162 registros novos. Os metadados excluíram 1.066, a triagem
por título e resumo excluiu 930, e 26 eram estudos secundários. Dos 140
relatos buscados, 43 não foram recuperados. Dos 97 avaliados no texto
completo, 63 foram excluídos, e 34 entraram. A RSL 1 tem, assim, 57 estudos
primários: 23 da busca nas bases e 34 do \textit{snowballing}.

% Gerado por docs/triagem/sintese/scripts/gera_tex.py. Não editar à mão.
\begin{table}[htbp]
\centering
\caption{Fluxo de seleção da RSL 1: busca nas bases}
\label{tab:prisma-rsl1}
\small
\begin{tabular}{p{2.2cm}p{9.5cm}r}
\hline
\textbf{Fase} & \textbf{Caixa} & \textbf{n} \\
\hline
Identificação & Registros das bases & 1.928 \\
 & Scopus: 882 + 65 da ampliação & 947 \\
 & Web of Science: 494 + 44 & 538 \\
 & IEEE Xplore: 266 exportados + 17 (eram 267 resultados) & 283 \\
 & ACM DL: 137 + 8 & 145 \\
 & Fontes nacionais (sem duplicados entre elas) & 15 \\
 & Removidos antes da triagem: duplicatas (786 pelo script, 3 conferidas no Rayyan) & 789 \\
Triagem & Registros triados (título, resumo e palavras-chave) & 1.139 \\
 & Excluídos na etapa 1 & 689 \\
 & Passaram para a etapa 2 & 450 \\
 & Excluídos pelo recorte, pelo resumo e metadados (CI6 187, CI5 88, CI7 86) & 361 \\
 & Estudos secundários, analisados à parte & 19 \\
 & Relatos buscados (texto completo) & 70 \\
 & Relatos não recuperados (CE4: UnB/CAPES e acesso aberto) & 21 \\
 & Relatos lidos (elegibilidade) & 49 \\
 & Excluídos no texto completo (CI5 13, CI6 6, CI1 4, CI2 2, CE1 1) & 26 \\
Inclusão & Estudos primários incluídos & 23 \\
\hline
\end{tabular}
\fonte{Elaborada pelo autor, no modelo de fluxograma do PRISMA 2020 \cite{page2021prisma}.}
\end{table}
\begin{table}[htbp]
\centering
\caption{Fluxo de seleção da RSL 1: \textit{snowballing}}
\label{tab:snowball-rsl1}
\small
\begin{tabular}{p{11.7cm}r}
\hline
\textbf{Caixa} & \textbf{n} \\
\hline
Registros novos identificados (referências e citações) & 2.162 \\
Excluídos pelos metadados & 1.066 \\
Triados por título e resumo & 1.096 \\
Excluídos na triagem & 930 \\
Estudos secundários (analisados à parte) & 26 \\
Relatos buscados & 140 \\
Relatos não recuperados & 43 \\
Relatos avaliados no texto completo & 97 \\
Excluídos no texto completo & 63 \\
\textbf{Estudos incluídos pelo \textit{snowballing}} & \textbf{34} \\
\hline
\end{tabular}
\fonte{Elaborada pelo autor, na coluna ``outros métodos do modelo do PRISMA 2020 \cite{page2021prisma}.}
\end{table}


Os relatos não recuperados (21 na busca nas bases e 43 no
\textit{snowballing}) são uma limitação desta revisão, e a impossibilidade de
acessar todos os relatos deve ser discutida \cite[item 23c]{page2021ee}. Os
estudos secundários (19 da busca nas bases e 26 do \textit{snowballing}) não
entram na síntese abaixo.
% PENDENTE: a lista dos relatos não recuperados está em
% docs/triagem/pedidos-autores.md. Registrar se os pedidos aos autores foram
% feitos e com que prazo (CE4).
% PENDENTE: os estudos secundários da RSL 1 (19 + 26) ainda não foram lidos.

% ---------------------------------------------------------------------------
\subsection{Características dos estudos}
\label{sec:rsl1-caracteristicas}

As Tabelas~\ref{tab:map-rsl1-anodepublica} a~\ref{tab:map-rsl1-fasedociclod}
resumem o mapeamento dos 57 estudos. Os estudos vão de 2015 a 2026, e 2023 e
2024 são os anos com mais estudos (9 cada). São 42 relatos em periódico (74\%)
e 15 em conferência (26\%). Os países mais frequentes do primeiro autor são
Espanha (8), Grécia (6) e Índia (6). Nenhum estudo é do Brasil ou de outro
país da América Latina.

A maioria dos estudos (45, ou 79\%) propõe uma ferramenta e a avalia. Nove
são estudos empíricos, e três têm outro tipo. Nas técnicas de IA (QP2),
predominam o PLN clássico (38 estudos, 67\%) e o aprendizado de máquina
clássico (36, 63\%). LLM e \textit{transformers} aparecem em 20 estudos
(35\%), e 5 estudos não usam IA ou não a relatam. Nas fases do ciclo de
inteligência, a IA e as ferramentas se concentram em processamento (53, 93\%),
análise (50, 88\%) e coleta (44, 77\%). Planejamento (6, 11\%) e avaliação
(3, 5\%) quase não aparecem.

% Gerado por docs/triagem/sintese/scripts/gera_tex.py. Não editar à mão.
\begin{table}[htbp]
\centering
\caption{RSL 1: ano de publicação dos estudos incluídos}
\label{tab:map-rsl1-anodepublica}
\small
\begin{tabular}{p{8cm}rr}
\hline
\textbf{Categoria} & \textbf{n} & \textbf{\%} \\
\hline
2015 & 2 & 4 \\
2016 & 1 & 2 \\
2017 & 1 & 2 \\
2018 & 3 & 5 \\
2019 & 8 & 14 \\
2020 & 7 & 12 \\
2021 & 2 & 4 \\
2022 & 1 & 2 \\
2023 & 9 & 16 \\
2024 & 9 & 16 \\
2025 & 8 & 14 \\
2026 & 6 & 11 \\
\hline
\end{tabular}
\fonte{Elaborada pelo autor.}
\end{table}
\begin{table}[htbp]
\centering
\caption{RSL 1: tipo de estudo dos estudos incluídos}
\label{tab:map-rsl1-tipodeestudo}
\small
\begin{tabular}{p{8cm}rr}
\hline
\textbf{Categoria} & \textbf{n} & \textbf{\%} \\
\hline
proposta de ferramenta com avaliação & 45 & 79 \\
estudo empírico & 9 & 16 \\
outro & 3 & 5 \\
\hline
\end{tabular}
\fonte{Elaborada pelo autor.}
\end{table}
\begin{table}[htbp]
\centering
\caption{RSL 1: família de técnica de ia dos estudos incluídos}
\label{tab:map-rsl1-famliadetcni}
\small
\begin{tabular}{p{8cm}rr}
\hline
\textbf{Categoria} & \textbf{n} & \textbf{\%} \\
\hline
PLN clássico (TF-IDF, NER, tópicos) & 38 & 67 \\
Aprendizado de máquina clássico & 36 & 63 \\
Aprendizado profundo (CNN, RNN, LSTM) & 28 & 49 \\
Grafos e análise de redes & 24 & 42 \\
LLM e transformers & 20 & 35 \\
Visão computacional e OCR & 14 & 25 \\
Agrupamento & 7 & 12 \\
sem IA ou não relatado & 5 & 9 \\
\hline
\end{tabular}
\fonte{Elaborada pelo autor. Um estudo pode ter mais de uma categoria.}
\end{table}
\begin{table}[htbp]
\centering
\caption{RSL 1: fase do ciclo de inteligência dos estudos incluídos}
\label{tab:map-rsl1-fasedociclod}
\small
\begin{tabular}{p{8cm}rr}
\hline
\textbf{Categoria} & \textbf{n} & \textbf{\%} \\
\hline
processamento & 53 & 93 \\
análise & 50 & 88 \\
coleta & 44 & 77 \\
difusão & 16 & 28 \\
planejamento & 6 & 11 \\
avaliação & 3 & 5 \\
\hline
\end{tabular}
\fonte{Elaborada pelo autor. Um estudo pode ter mais de uma categoria.}
\end{table}


A Tabela~\ref{tab:cobertura-rsl1} mostra quantos estudos trazem informação em
cada campo da extração. Os campos de ferramenta, fontes, ameaças e avaliação
estão preenchidos nos 57 estudos. Os campos mais vazios são o tratamento da
evidência (22 estudos, 39\%), o órgão usuário nomeado (25, 44\%) e os
requisitos legais e éticos (37, 65\%). Esses três vazios voltam como lacunas
nas QP3 e QP5.

% Gerado por docs/triagem/sintese/scripts/gera_tex.py. Não editar à mão.
\begin{table}[htbp]
\centering
\caption{RSL 1: estudos com informação em cada campo da extração}
\label{tab:cobertura-rsl1}
\small
\begin{tabular}{p{6cm}lrr}
\hline
\textbf{Campo} & \textbf{QP} & \textbf{n} & \textbf{\%} \\
\hline
Ferramenta e arquitetura & QP1 & 57 & 100 \\
Fontes OSINT & QP1 & 57 & 100 \\
Técnicas de IA & QP2 & 53 & 93 \\
Órgão usuário nomeado & -- & 25 & 44 \\
Requisitos legais e éticos & QP3 & 37 & 65 \\
Ameaças e riscos & QP4 & 57 & 100 \\
Método de avaliação & QP5 & 57 & 100 \\
Tratamento da evidência & QP5 & 22 & 39 \\
\hline
\end{tabular}
\fonte{Elaborada pelo autor.}
\end{table}


Os 32 estudos sem órgão usuário nomeado parecem contradizer o CI6, que exige
uso policial explícito. A diferença está no que cada um pede. O CI6 foi
aplicado ao registro e aceita a declaração de que a ferramenta é para a
polícia ou foi usada por ela. O campo da extração pede o nome de um órgão
concreto, e esses 32 estudos citam o uso policial de forma genérica.
% PENDENTE: R1 e R2 decidem se o CI6 exige órgão nomeado ou basta o uso
% policial declarado. A decisão vale para o grupo de 32 estudos, não só para
% RSL1-0856, cujo texto não nomeia nenhum órgão (qualidade/LEIAME.md).
% PENDENTE: o campo tipo_crime da extração ainda depende de R1 e R2 e não é
% relatado aqui.

Dois pares de estudos podem ser relatos do mesmo projeto: RSL1-0340 e
SB1-0202, ambos sobre o projeto SIGNIFICANCE, e SB1-0272 e SB1-0421, ambos
sobre detecção de atividades ilegais na \textit{dark web}. Até a decisão,
eles contam como estudos separados.
% PENDENTE: R1 e R2 conferem os pares RSL1-0340/SB1-0202 e SB1-0272/SB1-0421
% pelos autores. Se algum par for o mesmo estudo, a RSL 1 passa a ter 56 ou
% 55 estudos, e as tabelas são regeradas.

% ---------------------------------------------------------------------------
\subsection{Qualidade dos estudos}
\label{sec:rsl1-qualidade}

Na RSL 1, a nota média é 8,58 de 11, e a mediana é 9
(Tabela~\ref{tab:qualidade-resumo}). As notas vão de 4 a 11. Oito estudos
ficaram entre 4 e 6, 24 entre 7 e 9, e 25 entre 10 e 11. As notas mais baixas
são de RSL1-0856 (4), RSL1-0309 e RSL1-0368. A nota pondera a síntese e não
exclui estudos (Seção~\ref{sec:metodologia-qualidade}).

A Tabela~\ref{tab:qualidade-itens} mostra o número de ``sim'' em cada
critério. Todos os 57 estudos relatam pesquisa empírica (AQ1) e têm valor
para a pesquisa ou para a prática (AQ11). Os critérios com menos ``sim'' são
o AQ9 (2 estudos, 4\%), o AQ10 (34, 60\%), o AQ8 (39, 68\%) e o AQ6 (40,
70\%).

O AQ9 pergunta se a relação entre o pesquisador e os participantes foi
considerada. Só RSL1-0381 e SB1-0781 têm ``sim''. Pela regra de aplicação
desse critério, conta como ``sim'' discutir como os pesquisadores podem ter
influenciado os resultados, por exemplo ao avaliar a própria ferramenta, e,
em estudos sem participantes, discutir o viés dos autores como avaliadores.
Quase nenhum estudo faz essa discussão. Como 45 dos 57 estudos propõem uma
ferramenta e a avaliam, o resultado do AQ9 indica que a maior parte da
evidência da RSL 1 vem de avaliações feitas pelos autores da ferramenta, sem
discussão desse viés. Esse achado reforça a fragilidade da
validade externa descrita na QP5 (Seção~\ref{sec:rsl1-qp5}). Entre os
motivos de ``não'' no AQ10 (achados críveis) estão números inconsistentes no
próprio relato, como em RSL1-0340 e SB1-0421.
% PENDENTE: o "sim" de SB1-0781 no AQ9 é limítrofe (discussão de suposições
% de rotulagem, não de viés pessoal; qualidade/rsl1-qualidade.csv). R1 e R2
% confirmam.
% PENDENTE: as notas são da IA (terceira opinião). O protocolo pede avaliação
% independente de R1 e R2, com kappa, ainda não feita.

% ---------------------------------------------------------------------------
\subsection{QP1 e QP2: ferramentas, arquiteturas e técnicas de IA}
\label{sec:rsl1-qp1-qp2}

A síntese temática seguiu as três etapas de \citeonline{thomas2008}:
codificação, temas descritivos e temas analíticos. A
Tabela~\ref{tab:temas-A1} lista os temas analíticos das QP1 e QP2. Nessa
tabela e nas seguintes, a coluna ``Estudos'' conta os estudos distintos nos
códigos do tema, inclusive os códigos que registram uma ausência (por
exemplo, ``nenhum mecanismo de preservação descrito''). O número mostra a
abrangência do tema no corpus, e não quantos estudos sustentam a afirmação do
tema.

% Gerado por docs/triagem/sintese/scripts/gera_tex.py. Não editar à mão.
\begin{table}[htbp]
\centering
\caption{Temas analíticos: RSL 1, QP1 e QP2}
\label{tab:temas-A1}
\small
\begin{tabular}{lp{11cm}r}
\hline
\textbf{Tema} & \textbf{Descrição} & \textbf{Estudos} \\
\hline
TA1 & Arquitetura de referência: pipeline modular com busca do investigador no centro & 56 \\
TA2 & Funções de IA: triagem, extração e ligação, com LLM como camada de apoio & 52 \\
TA3 & Distância entre modelos de pesquisa e ferramentas adotadas pela polícia & 57 \\
\hline
\end{tabular}
\fonte{Elaborada pelo autor. ``Estudos: número de estudos distintos nos códigos do tema.}
\end{table}


\textbf{TA1. Arquitetura de referência: \textit{pipeline} modular, com a
busca do investigador no centro.} Os estudos convergem num fluxo de coleta,
normalização, análise e apresentação (RSL1-0288, SB1-0421, SB1-0837). Os
sistemas mais completos separam conectores por fonte, armazenamento indexado
e módulos de análise, e expõem tudo por API e painel (SB1-0202, SB1-0421,
RSL1-0856, SB1-0837). A busca ativa pelo investigador é rara: o código
correspondente reúne 8 estudos, entre eles RSL1-0183, SB1-0444, RSL1-0494 e
RSL1-0856. A maioria entrega só um classificador ou componente, sem
interface (RSL1-0211, SB1-0767, SB1-0914). As fontes definem os conectores.
A \textit{dark web} é fonte em 33 estudos e exige \textit{crawler} para Tor,
com sessão e CAPTCHA (RSL1-0309, SB1-0629, SB1-0436). Redes sociais,
notícias, infraestrutura técnica e \textit{blockchain} exigem APIs e
raspadores próprios (SB1-0404, RSL1-0381, RSL1-0908). Muitos estudos validam
em conjuntos de dados prontos, e não em coleta viva (SB1-0367, SB1-0547,
SB1-0873).

\textbf{TA2. Funções de IA: triagem, extração e ligação, com LLM como camada
de apoio.} O uso central da IA é triar conteúdo, isto é, classificar textos e
imagens como ilícitos ou relevantes (SB1-0620, SB1-0547, SB1-0827, SB1-0837).
Depois vêm a extração de entidades e campos (SB1-0724, SB1-0873), a ligação
de contas, \textit{aliases} e redes (RSL1-0466, SB1-0781, SB1-0767) e a
localização no tempo e no espaço (RSL1-0183, RSL1-0662). Os LLMs fazem
extração, modelagem de tópicos, checagem e busca semântica (SB1-0724,
SB1-1001, RSL1-0381, RSL1-0494), em geral seguidos de expressões regulares
ou de verificação humana (SB1-0724, SB1-0837). Sumarização e expansão de
consultas apoiam o analista (RSL1-0494, SB1-1025, SB1-0444). Técnicas para
poucos rótulos mostram a falta de dados rotulados na área (SB1-0678,
SB1-0828, SB1-0367). A detecção de conteúdo sintético é rara (RSL1-0381,
RSL1-0466).

\textbf{TA3. Distância entre os modelos de pesquisa e as ferramentas
adotadas pela polícia.} As ferramentas que a polícia de fato usa são
navegadores sem IA, centrados em auditoria, captura, \textit{hash} e
relatório (RSL1-0854, RSL1-0863). Processos reais combinam ferramentas de
terceiros (RSL1-0908, RSL1-0466, RSL1-0816). Já a maioria das propostas com
IA não tem órgão usuário concreto (SB1-0298, SB1-0547, SB1-0620). A IA fica
em processamento e análise, e planejamento e avaliação quase não recebem
apoio (RSL1-0214, RSL1-0368, RSL1-0612). Poucos estudos unem IA e
preservação de evidência no mesmo sistema (SB1-0306, RSL1-0466, SB1-0837), e
a explicabilidade é rara (SB1-0781, SB1-0837). Os modelos de avaliação de
ferramentas mostram que o órgão julga segurança, facilidade de uso e
acurácia, e não só o desempenho do modelo (RSL1-0612, RSL1-0381).

\textbf{Contradições.} Primeiro, as ferramentas adotadas por polícias não
usam IA e priorizam evidência e relatório (RSL1-0854, RSL1-0863), enquanto a
maioria das propostas se centra em classificadores sem órgão usuário
(SB1-0547, SB1-0620); o modelo de valor do Exército dos EUA trata o
aprendizado de máquina como tecnologia futura (RSL1-0612). Segundo, a
autonomia da decisão diverge: um sistema emite alerta automático de conteúdo
ilegal (SB1-0465), e outros exigem decisão final humana (RSL1-0368,
SB1-0837). Terceiro, técnicas tradicionais seguem como alternativa às
recentes, como SIFT com SVM diante de \textit{features} de CNN (SB1-0827) e
LDA e TF-IDF diante de LLM (SB1-1001). Quarto, o projeto SIGNIFICANCE
aparece com arquiteturas diferentes em RSL1-0340 e SB1-0202.
% PENDENTE: se R1 e R2 decidirem que RSL1-0340 e SB1-0202 são o mesmo
% estudo, a quarta contradição vira evolução do projeto entre dois relatos.

\textbf{Lacunas.} Nenhum estudo relata órgão usuário, fonte ou idioma do
Brasil ou da América Latina. Planejamento e avaliação do ciclo quase não
recebem apoio de ferramenta. Mensageiros e grupos fechados são pouco
cobertos como fonte (RSL1-0368, RSL1-0466, RSL1-0816, SB1-0689, SB1-0837).
IA e preservação de evidência raramente aparecem juntas (SB1-0306,
RSL1-0466, SB1-0837, SB1-0421). A explicabilidade quase não aparece
(SB1-0781, SB1-0837), e em RSL1-0309 fica como trabalho futuro.

% ---------------------------------------------------------------------------
\subsection{QP3 e QP4: requisitos legais e éticos, ameaças e mitigações}
\label{sec:rsl1-qp3-qp4}

A Tabela~\ref{tab:temas-A2} lista os temas analíticos das QP3 e QP4. O
recorte da etapa 2 deixou de fora estudos conceituais, legais e éticos sem
avaliação de ferramenta (CI5). Por isso, a evidência das QP3 e QP4 vem
sobretudo do que os estudos de ferramenta relatam sobre esses temas, e 20
dos 57 estudos não tratam de requisitos legais nem éticos.

% Gerado por docs/triagem/sintese/scripts/gera_tex.py. Não editar à mão.
\begin{table}[htbp]
\centering
\caption{Temas analíticos: RSL 1, QP3 e QP4}
\label{tab:temas-A2}
\small
\begin{tabular}{lp{11cm}r}
\hline
\textbf{Tema} & \textbf{Descrição} & \textbf{Estudos} \\
\hline
TA1 & Conformidade legal é requisito de concepção, não etapa posterior & 36 \\
TA2 & Evidência só tem valor com proveniência e integridade verificáveis & 28 \\
TA3 & IA gera pistas; a decisão e a explicação ficam com o humano & 35 \\
TA4 & O modelo de ameaças tem três alvos: investigador, ferramenta e dados & 22 \\
TA5 & A coleta em ambiente adversarial exige resiliência e manutenção contínua & 49 \\
\hline
\end{tabular}
\fonte{Elaborada pelo autor. ``Estudos: número de estudos distintos nos códigos do tema.}
\end{table}


\textbf{TA1. A conformidade legal é requisito de concepção, e não etapa
posterior.} Os estudos que discutem o direito tratam a conformidade como
restrição de projeto. Na polícia holandesa, a leitura do GDPR proíbe
identificar usuários, ligar perfis, entrar em grupos fechados e armazenar
dados brutos (RSL1-0183). No Reino Unido, o nível da investigação define o
que é permitido, sob necessidade e proporcionalidade (RSL1-0854). Os testes
passam por avaliação de impacto e minimização (RSL1-0381, SB1-0781), e as
ações ativas ou intrusivas dependem de autorização (RSL1-0368, RSL1-0816,
RSL1-0908). A jurisdição limita o alcance (RSL1-0561, SB1-0689). Na coleta
automatizada não há consenso, e parte dos estudos contorna barreiras sem
discutir a legalidade (SB1-0629, SB1-0436, SB1-0678).

\textbf{TA2. A evidência só tem valor com proveniência e integridade
verificáveis.} A admissibilidade em juízo depende da proveniência e da
integridade (RSL1-0466, RSL1-0854, SB1-0104, SB1-0528, SB1-0837). Como o
conteúdo é volátil, a captura precisa ser imediata: publicações são editadas
ou apagadas (RSL1-0466), e a página pode sumir entre a anotação e o
arquivamento (SB1-0528). As respostas técnicas são o \textit{hash}
(SB1-0306), o \textit{blockchain} ligado a normas ISO/IEC e às regras de
prova dos EUA (RSL1-0466) e a trilha de auditoria automática, porque o
registro manual falha (RSL1-0863).

\textbf{TA3. A IA gera pistas, e a decisão e a explicação ficam com o
humano.} Os estudos citam como riscos da IA o erro de classificação
(RSL1-0348, RSL1-0618), a alucinação de LLM (RSL1-0494, RSL1-0906,
RSL1-0908, SB1-0724), o viés (RSL1-0368, SB1-0584, SB1-1001), a opacidade
(RSL1-0381, SB1-0367, SB1-0781) e o viés de automação (RSL1-0381,
RSL1-0618). A resposta mais frequente trata os vereditos como pistas, que
exigem corroboração independente e revisão humana (RSL1-0381). Ela também
ancora a resposta no texto da fonte (RSL1-0494) e verifica a saída com
técnicas de explicação (SB1-0837). Várias mitigações de viés e de
explicabilidade ficam como trabalho futuro (SB1-0367, SB1-0547).

\textbf{TA4. O modelo de ameaças tem três alvos: investigador, ferramenta e
dados.} O investigador deixa rastros nos registros dos alvos (RSL1-0381),
expõe IP e localização (RSL1-0854) e pode alertar o suspeito (RSL1-0908,
RSL1-0816), e os sites detectam o coletor (SB1-0421, SB1-0528, SB1-1026). Na
ferramenta, o risco de ataque justifica o sigilo (RSL1-0183), e uma
infraestrutura frágil compromete a segurança (RSL1-0906, SB1-0528). Nos
dados, pesam a transferência a terceiros, como \textit{chatbots} externos
(RSL1-0906), e a dependência de \textit{software} estrangeiro (RSL1-0856).
As mitigações relatadas são o anonimato do operador e ambientes virtuais
(RSL1-0381, RSL1-0309, SB1-0272), os dados mantidos localmente (RSL1-0854)
e a infraestrutura controlada pelo órgão (RSL1-0856).

\textbf{TA5. A coleta em ambiente adversarial exige resiliência e manutenção
contínua.} Os sites somem e migram (SB1-0767, SB1-0877, SB1-0914), e as
barreiras contra a coleta se multiplicam (SB1-0436, SB1-1026). Os criminosos
ofuscam termos e trocam de identidade (SB1-0272, SB1-0689, SB1-0781). As
gírias mudam e degradam modelos treinados numa base estática (SB1-0367,
SB1-0873). As ferramentas ficam obsoletas (RSL1-0906), e as plataformas
mudam termos de uso e APIs (RSL1-0381, SB1-0584). Na coleta, as mitigações
são novas tentativas, espera crescente e redescoberta de endereços
(SB1-0306, RSL1-0309). Nos modelos, o aprendizado contínuo aparece só como
proposta (SB1-0367).

\textbf{Contradições.} Seis pontos dividem os estudos. Sobre o
\textit{robots.txt}, RSL1-0309 e SB1-0104 respeitam o arquivo, SB1-0629 o
ignora, e SB1-0421 oferece a opção de ignorá-lo. A ligação de identidades é
proibida na leitura holandesa do GDPR (RSL1-0183), mas RSL1-0466 correlaciona
conta e pessoa, e SB1-0781 liga \textit{aliases} sob avaliação de impacto. A
entrada em grupos fechados é proibida em RSL1-0183, admitida com autorização
em RSL1-0368 e RSL1-0612 e feita sem discussão ética em SB1-0689. A recusa
ética dos LLMs é lacuna a contornar em SB1-0724 e salvaguarda em SB1-0367. A
captura automática é criticada como excessiva em RSL1-0854, enquanto
RSL1-0863 mostra que o registro manual falha. Por fim, RSL1-0863 apresenta o
código aberto como mitigação de confiança, e RSL1-0183 relata o sigilo das
ferramentas por risco de ataque.

\textbf{Lacunas.} Dez estudos monitoram pessoas, grupos ou manifestações sem
discussão substantiva de ética ou da privacidade dos monitorados: RSL1-0662,
RSL1-0288, SB1-0867, RSL1-0344, SB1-0297, SB1-0689, SB1-0421, RSL1-0561,
RSL1-0348 e SB1-0404. Nos dois últimos, o caso é limítrofe. Nenhum estudo
trata da lei brasileira. A soberania dos dados aparece só em RSL1-0856,
RSL1-0854, RSL1-0494 e RSL1-0906. Ameaças contra a própria ferramenta
aparecem só em RSL1-0183, RSL1-0856, RSL1-0863, RSL1-0906 e SB1-0528, e
nenhum estudo relata teste de segurança. Quatro estudos contornam barreiras
de acesso sem discussão legal ou ética (SB1-0436, SB1-0629, SB1-0678,
SB1-0877).
% PENDENTE: R1 e R2 confirmam a classificação de RSL1-0348 e SB1-0404 como
% casos limítrofes do grupo de vigilância sem discussão ética.

% ---------------------------------------------------------------------------
\subsection{QP5: avaliação das ferramentas e preservação da evidência}
\label{sec:rsl1-qp5}

A Tabela~\ref{tab:temas-A3} lista os temas analíticos da QP5. Em 38 dos 57
estudos, a avaliação dominante são métricas num conjunto de dados. Só 7
estudos avaliam com usuários. Só 7 tratam explicitamente da cadeia de
custódia, e outros 15 descrevem algum registro ou preservação, como
\textit{hash}, carimbo de tempo ou captura de tela, sem tratar da cadeia de
custódia. Nenhum estudo avalia se a evidência é aceita em juízo.

% Gerado por docs/triagem/sintese/scripts/gera_tex.py. Não editar à mão.
\begin{table}[htbp]
\centering
\caption{Temas analíticos: RSL 1, QP5}
\label{tab:temas-A3}
\small
\begin{tabular}{lp{11cm}r}
\hline
\textbf{Tema} & \textbf{Descrição} & \textbf{Estudos} \\
\hline
TA1 & A avaliação mede o modelo, quase nunca o trabalho do investigador & 56 \\
TA2 & A validade externa das avaliações é frágil & 42 \\
TA3 & Coleta, cobertura e desempenho operacional são critérios de avaliação próprios & 54 \\
TA4 & Preservação da evidência é exceção, mas quem a trata converge para um mesmo conjunto de mecanismos & 57 \\
TA5 & A saída da IA orienta a investigação, mas não substitui a prova & 56 \\
\hline
\end{tabular}
\fonte{Elaborada pelo autor. ``Estudos: número de estudos distintos nos códigos do tema.}
\end{table}


\textbf{TA1. A avaliação mede o modelo e quase nunca o trabalho do
investigador.} Os estudos avaliam o componente de IA isolado, com métricas de
classificação em dados separados para teste ou em validação cruzada, e
comparam com \textit{baselines} (RSL1-0561, SB1-0367, SB1-0584). Em outras
tarefas, usam métricas de agrupamento, tópicos, ordenação ou resumo
(RSL1-0344, RSL1-0618, SB1-0914, SB1-1025). Os 7 estudos com usuários são
RSL1-0183, RSL1-0381, RSL1-0494, RSL1-0854, RSL1-0863, RSL1-0906 e SB1-0837.
Em RSL1-0494, o perfil dos participantes não está no campo extraído. Os
resultados com participantes divergem. Em RSL1-0494, o grupo com IA teve
pontuação maior (18,2 contra 11,5; p = 0,007). Em RSL1-0183, as hipóteses
H1 a H4 não foram confirmadas, mas elas testam outra coisa. O SUS
(RSL1-0183, RSL1-0494) e o tempo de tarefa (RSL1-0906, RSL1-0368) têm
direção favorável à ferramenta, e parte do benefício só é declarada, sem
dados (RSL1-0863, SB1-0202).

\textbf{TA2. A validade externa das avaliações é frágil.} A verdade de
referência vem de bases públicas, de rótulos por heurística (RSL1-0288,
SB1-0668) ou de anotação manual pequena, como as 30 páginas de SB1-0724. A
validação contra dado independente é rara (RSL1-0779, SB1-0867, SB1-1026).
Estudos de caso sem métricas (RSL1-0466, RSL1-0908) e validação julgada
pelos próprios autores (RSL1-0214) não permitem comparação. Há dados ou
denominadores não descritos (RSL1-0381, RSL1-0856, SB1-0465) e números
inconsistentes (RSL1-0340, SB1-0421). Quando policiais testam ferramentas
existentes, o desempenho varia muito, de 100\% na busca reversa de imagens a
16\% numa ferramenta de notícias (RSL1-0381).

\textbf{TA3. Coleta, cobertura e desempenho operacional são critérios de
avaliação próprios.} Os estudos de coleta medem vazão, latência, taxa de
sucesso e cobertura (RSL1-0309, SB1-0421, SB1-0104, SB1-0629). Cobertura de
sites e cobertura de atividade são medidas diferentes: em SB1-1026, a coleta
cobre 44,6\% dos sites e 98\% da atividade. O registro de integridade tem
custo que cresce com o volume (RSL1-0466). Componentes centrais ficam sem
avaliação, como o \textit{crawler} (SB1-0436) e os extratores (RSL1-0309).

\textbf{TA4. A preservação da evidência é exceção, e quem a trata converge
para os mesmos mecanismos.} Muitos estudos não descrevem preservação, ou
declaram que não tratam da cadeia de custódia (RSL1-0348, SB1-0724). Os 7
estudos que a tratam são RSL1-0466, RSL1-0816, RSL1-0854, RSL1-0863,
SB1-0306, SB1-0528 e SB1-0837. Os mecanismos se repetem: \textit{hash} na
captura, carimbo de tempo, conteúdo bruto, registro das ações, proveniência
e relatório exportável (RSL1-0466, RSL1-0854, RSL1-0863, SB1-0528,
SB1-0104). A exigência vem do contexto jurídico, como o precedente Bucknor
v R, citado em RSL1-0854, e a trilha de auditoria das diretrizes da ACPO,
citada em RSL1-0863. Selagem e \textit{ledger} aparecem como
reforço (RSL1-0466, SB1-0306, SB1-0421), mas em SB1-0306 o \textit{ledger} é
simulado.
% PENDENTE: R1 e R2 conferem a classificação de cadeia de custódia (A3,
% "Decisões de codificação"): SB1-0837 ("estilo cadeia de custódia") e
% SB1-0306 (ledger simulado) marcados "sim"; SB1-0104 e SB1-0421 marcados
% "não".

\textbf{TA5. A saída da IA orienta a investigação, mas não substitui a
prova.} RSL1-0381 e SB1-0781 afirmam que os resultados automáticos não
sustentam sozinhos uma alegação probatória. Outros estudos buscam
rastreabilidade: RSL1-0494 liga as respostas ao texto da fonte, e SB1-0272
escolhe a fusão linear de modelos pela interpretabilidade. RSL1-0618
discute a confiabilidade de processos automatizados em tribunal, e RSL1-0816
sustenta a validade com um procedimento controlado e repetível. Já
RSL1-0381 não preserva os registros de execução.

\textbf{Contradições.} Sobre reter dados, RSL1-0183 não guarda dados brutos,
e RSL1-0466, RSL1-0854, RSL1-0863 e SB1-0528 guardam o conteúdo bruto com
\textit{hash} e registro. O desempenho dos LLMs é instável: em SB1-0367, o
SVM supera o Llama 3.2 na tarefa binária e perde na multiclasse, e em
SB1-0837 o GPT-4o mini fica abaixo do modelo proposto. Em RSL1-0618, a
métrica automática piora, mas os tópicos ficam mais coerentes à inspeção
humana.

\textbf{Lacunas.} Fora dos 7 estudos com usuários e dos 7 com cadeia de
custódia citados acima, nenhum estudo avalia com usuários ou trata da cadeia
de custódia. Nenhum estudo avalia a aceitação da evidência em juízo, e
RSL1-0466 registra essa ausência. O uso operacional prolongado só aparece no
OSIRT, um navegador para investigadores (RSL1-0854, RSL1-0863).
Componentes centrais ficam sem avaliação em SB1-0436, RSL1-0309, RSL1-0466 e
SB1-0629. Nenhum estudo usa dados em português.

% ---------------------------------------------------------------------------
\subsection{Rastreabilidade para o buscador}
\label{sec:rsl1-rastreabilidade}

A última leitura cruzou os achados com o levantamento de requisitos do
buscador (Seção~\ref{sec:metodologia-requisitos}). Ela reuniu os requisitos
candidatos dos estudos em 27 temas descritivos, ligados a um requisito do
levantamento ou marcados como novos (Tabela~\ref{tab:rastreabilidade}), e em
7 temas analíticos (Tabela~\ref{tab:temas-A4}).

% Gerado por docs/triagem/sintese/scripts/gera_tex.py. Não editar à mão.
\begin{table}[htbp]
\centering
\caption{Temas analíticos: RSL 1, requisitos}
\label{tab:temas-A4}
\small
\begin{tabular}{lp{11cm}r}
\hline
\textbf{Tema} & \textbf{Descrição} & \textbf{Estudos} \\
\hline
TA1 & A revisão confirma o núcleo probatório: custódia, auditoria e relatório & 21 \\
TA2 & A revisão confirma o grafo e pede cautela na atribuição & 21 \\
TA3 & A revisão amplia o RNF05: IA explicada, como pista, com revisão humana & 39 \\
TA4 & A revisão confirma e detalha os requisitos legais e operacionais & 23 \\
TA5 & A revisão confirma a coleta multifonte e o monitoramento, mas em outras fontes & 46 \\
TA6 & Requisitos do domínio de corrupção e lavagem ficam sem respaldo & 45 \\
TA7 & Usabilidade, desempenho e gestão de casos têm apoio modesto & 23 \\
\hline
\end{tabular}
\fonte{Elaborada pelo autor. ``Estudos: número de estudos distintos nos códigos do tema.}
\end{table}


% Gerado por docs/triagem/sintese/scripts/gera_tex.py. Não editar à mão.
\begin{small}
\begin{longtable}{lp{6.8cm}llr}
\caption{RSL 1: requisitos consolidados e ligação com o levantamento de requisitos}\rotulounico{tab:rastreabilidade}\\
\hline
\textbf{ID} & \textbf{Requisito consolidado} & \textbf{Requisito} & \textbf{QP} & \textbf{Estudos} \\ \hline\endhead
TD1 & Coleta multifonte em fontes abertas & RF01 & QP1 & 22 \\
TD2 & Coleta na dark web & novo & QP1 & 22 \\
TD3 & Monitoramento contínuo e alertas & RF06 & QP1 & 15 \\
TD4 & Resiliência e controle da coleta & novo & QP1, QP4 & 14 \\
TD5 & Cadeia de custódia na captura & RF07 & QP5 & 16 \\
TD6 & Trilha de auditoria das ações & RNF02 & QP4, QP5 & 7 \\
TD7 & Relatório exportável como prova & RF09 & QP5 & 4 \\
TD8 & Interoperabilidade & RNF09 & QP1 & 4 \\
TD9 & Resolução de entidades e pivô por identificadores & RF02 & QP1, QP2 & 7 \\
TD10 & Grafo de vínculos & RF03 & QP1, QP2 & 9 \\
TD11 & Priorização de alvos e equilíbrio de erros & novo & QP2 & 9 \\
TD12 & Classificação automática de conteúdo & novo & QP2 & 25 \\
TD13 & Extração de entidades e uso de LLM & novo & QP2 & 15 \\
TD14 & Busca, filtros e triagem & novo & QP1, QP2 & 10 \\
TD15 & Análise geográfica e temporal & novo & QP2 & 9 \\
TD16 & Multilíngue e linguagem local & novo & QP2 & 9 \\
TD17 & Explicabilidade e confiança das saídas de IA & RNF05 & QP2, QP5 & 16 \\
TD18 & IA como pista, revisão humana e confiabilidade das fontes & novo & QP2, QP3, QP5 & 13 \\
TD19 & Anonimato operacional & RNF03 & QP4 & 11 \\
TD20 & Privacidade e minimização & RNF04 & QP3 & 12 \\
TD21 & Autorização legal e conformidade processual & RNF04 & QP3 & 7 \\
TD22 & Segurança e soberania da infraestrutura & RNF01 & QP3, QP4 & 9 \\
TD23 & Usabilidade para investigadores & RNF08 & QP5 & 6 \\
TD24 & Desempenho e escala & RNF06 & QP1 & 6 \\
TD25 & Manutenção, código aberto e dependência de plataformas & RNF10 & QP1, QP4 & 6 \\
TD26 & Gestão de casos & RF08 & QP1 & 4 \\
TD27 & Disponibilidade e redundância & RNF07 & QP4 & 1 \\
\hline
\multicolumn{5}{p{13cm}}{\footnotesize Fonte: elaborada pelo autor. ``Requisito: requisito do levantamento inicial que já cobre o tema, ou ``novo.}\
\end{longtable}
\end{small}


\textbf{O que a revisão confirma.} A revisão confirma o núcleo probatório.
Para o orientador, o requisito mais crítico é a combinação de cadeia de
custódia (RF07) com trilha de auditoria (RNF02), e é essa combinação que os
estudos mais ligam à admissibilidade. \textit{Hash}, carimbo de tempo e
captura integral (RSL1-0466, RSL1-0863, SB1-0528, SB1-0837) andam junto com
o registro automático das ações (RSL1-0854, RSL1-0863) e com o relatório
exportável, que apoia o RF09 (RSL1-0854, RSL1-0466). Os grafos de vínculos
apoiam o RF03 (RSL1-0344, RSL1-0368, RSL1-0856, RSL1-0908, SB1-0867). O RF02
tem apoio indireto, porque os estudos ligam contas e \textit{aliases}, e não
CPF ou CNPJ. A coleta multifonte e o monitoramento com alerta apoiam o RF01 e
o RF06 (RSL1-0348, RSL1-0856, SB1-0837, RSL1-0908, SB1-0421, SB1-0877). O
anonimato operacional apoia o RNF03 (RSL1-0309, SB1-0421, RSL1-0854), e a
autorização para a coleta ativa e o registro das justificativas apoiam o
RNF04 (RSL1-0908, RSL1-0816, RSL1-0854, RSL1-0612, RSL1-0368). Usabilidade e
gestão de casos têm apoio menor, em poucos estudos (RSL1-0612, SB1-0528,
RSL1-0908, RSL1-0214).

\textbf{O que a revisão amplia.} O RF07 ganha quatro pontos: identificar o
coletor e a linha de base da coleta (RSL1-0466, RSL1-0816), verificar a
integridade depois da captura (SB1-0528), versionar e comparar coletas
repetidas (SB1-0421, RSL1-0816) e registrar a versão e a configuração do
modelo de IA (RSL1-0381). O RNF05 deixa de ser só explicar o alerta. A saída
da IA é pista e não conclusão (RSL1-0381), a decisão final é humana
(RSL1-0368, SB1-0202), e a saída do LLM é validada por esquema (SB1-0837,
SB1-0724). O RNF04 ganha a minimização por finalidade: capturar só o
necessário e não guardar dados pessoais sem necessidade (RSL1-0183,
RSL1-0854, SB1-1001). O RNF10 ganha a resiliência da coleta, porque as
fontes somem, mudam ou bloqueiam o coletor (SB1-0767, SB1-0306, RSL1-0381).
A revisão também sugere requisitos que o levantamento não tem: busca e
filtros, classificação de conteúdo, extração de entidades com LLM e fonte
citada, análise geográfica e temporal, suporte multilíngue, resiliência da
coleta e priorização de alvos com uma política de erro declarada, isto é, a
escolha entre precisão e revocação (RSL1-0561, SB1-0584).

\textbf{O que fica sem respaldo.} Nenhum estudo pede as tipologias de
corrupção (RF04), a linha do tempo patrimonial (RF05), as integrações sob
ordem judicial (RF10) ou o DevSecOps (RNF11). Também não aparecem as fontes
oficiais brasileiras do RF01, a resolução por CPF ou CNPJ do RF02, a
autenticação multifator, o acesso por perfil e a criptografia do RNF01, nem a
exportação para ferramentas de análise de vínculos do RNF09. O RNF07 tem
apoio só em RSL1-0816. Os estudos tratam de terrorismo, armas, drogas, bens
culturais, cibercrime e mercados ilícitos (RSL1-0348, RSL1-0211, RSL1-0779,
RSL1-0340, SB1-0547), e o mais próximo da investigação financeira é o
rastreamento de criptoativos (RSL1-0908). A revisão sustenta, portanto, os
requisitos transversais do buscador, e não os requisitos do domínio de
corrupção e lavagem de dinheiro. Esses requisitos precisam ser justificados
pelo levantamento com o especialista e pela legislação. A falta de
ferramentas OSINT voltadas a corrupção é uma lacuna da literatura revisada.

\textbf{Contradições que afetam o produto.} Duas tensões pesam no projeto.
A primeira opõe minimização e preservação: RSL1-0183 não guarda dados
brutos, SB1-1001 não guarda dados pessoais, e RSL1-0854 critica a captura
automática de todas as páginas, enquanto SB1-0837, SB1-0724, SB1-0421 e
SB1-0528 guardam a cópia bruta para a proveniência. A segunda separa o
artefato capturado da inferência da IA: SB1-0306 e RSL1-0466 buscam a
admissibilidade do que foi coletado, e SB1-0781 e RSL1-0381 desaconselham
usar a saída analítica como prova.

\textbf{Implicações para o protótipo.} Os achados sustentam manter no
protótipo a cadeia de custódia na captura (RF07) e a trilha de auditoria
(RNF02), e acrescentar aos metadados de cada captura o coletor e a versão do
conector. A coleta fica passiva, com justificativa obrigatória para cada
busca. A ligação entre entidades fica como sugestão, com confirmação humana
quando só o nome coincide. Todo alerta gerado por IA é tratado como pista,
mostra o motivo e a fonte e exige a confirmação do analista, registrada na
auditoria. Os conectores registram falhas e mudanças de formato da fonte. A
captura deve ser seletiva e decidida pelo investigador, o que concilia a
integridade da evidência com a minimização. O monitoramento contínuo (RF06),
hoje fora do recorte do protótipo, tem respaldo na revisão. A avaliação do
protótipo deve ter dois níveis, como sugere a QP5: métricas por componente
num conjunto de teste próprio e documentado, e avaliação com policiais, com
usabilidade e tempo de tarefa comparado ao fluxo manual.
% PENDENTE: R1 e R2 decidem quais requisitos novos sugeridos pela leitura A4
% entram em docs/requisitos-app.md. Nada foi alterado lá.
% PENDENTE: a tensão entre minimização e preservação bruta afeta o RF07 do
% protótipo, que guarda a carga bruta. Decidir a política de retenção.
